\section{Tightness at the two smallest nonabelian orders}
\begin{theorem}[Uniform small-order reversal]\label{thm:small}
Let $G=\gen{a,b}$ have order $2p^m$, where $m=3$ or $4$.
Assume $|a|,|b|\geq3$, $(ab)^2=1$,
$\gen a\cap\gen b=1$ and $|a||b|<|G|$.
Then $G$ has an automorphism simultaneously inverting $a,b$.
In particular every chiral polyhedron of either order is tight.
\end{theorem}
\begin{proof}
The Sylow subgroup $P$ of odd order is normal of index two.
Exactly one of $a,b$ is in $P$: neither both inside (generation)
nor both outside (their nonidentity involution product would be
inside). Exchange them if necessary and put $x=a$, $y=b^2$.
Conjugation by $b$ satisfies $x^b=y^{-1}x^{-1}$ and $y^b=y$,
so $P=\gen{x,y}$. Their cyclic subgroups are nontrivial and disjoint.
If $P$ were abelian, its independent cyclic coordinates give an
inverting automorphism of the presentation
\[
\gen{a,b\mid a^{|a|},\ b^{|b|},\ (ab)^2,\ [a,b^2]}:
\]
collection bounds its order by $|a||b|$, contradicting non-tightness.
Thus $P$ is nonabelian. Write $|x|=p^u$, $|y|=p^v$.
Disjointness and non-tightness give $u,v\geq1$, $u+v<m$.

For $m=3$ we have $u=v=1$ and class two. With $c=[x,y]$,
the relations $x^p=y^p=c^p=[c,x]=[c,y]=1$ collect words into
$x^iy^jc^k$, at most $p^3$ possibilities. Adjoin $b$ with
$b^2=y$ and $(xb)^2=1$. The subgroup generated by $x,b^2$
is normal by the conjugation formula and has index at most two.
The resulting presented group maps onto $G$, and the bound $2p^3$
shows it is exactly $G$. Simultaneous inversion of $a,b$ respects
all relations: it inverts $x,y$, fixes their commutator up to
conjugation, and preserves class two. The resulting endomorphism
squares to the identity and is the required automorphism.

Let $m=4$. The possible pairs $(u,v)$ are $(1,1),(1,2),(2,1)$.
If the class is two, $P'$ is cyclic, generated by $[x,y]$;
the small-kernel lemma of~\cite{UF} therefore gives $|P'|=p$.
The pair $(1,1)$ would bound the order by $p^3$ and is impossible.
For the other two pairs, the relations
\[
x^{p^u}=y^{p^v}=c^p=[c,x]=[c,y]=1
\]
collect into at most $p^{u+v+1}=p^4$ words. Exactly the preceding
presentation and reversal argument applies.

Suppose the class is three. Put $c=[x,y]$, $d=[c,x]$,
$e=[c,y]$. Impose all weight-four commutators equal to one,
and $c^p=d^p=e^p=1$. These are actual relations by
the small-kernel lemma of~\cite{UF}. If $u=v=1$, the two elements $d,e$
generate the cyclic group $\gamma_3(P)$, so there is a relation
$d^i e^j=1$ with $(i,j)\ne(0,0)$ modulo $p$.
The corresponding presentation of $P$ has at most $p^4$ collected
words $x^\alpha y^\beta c^\gamma d^\delta e^\epsilon$:
the last two exponents have only $p$ independent possibilities.
If $(u,v)=(2,1)$, then $x^p$ is central, since
$[x^p,y]=c^p d^{\binom p2}=1$.
It is nontrivial, so generates $Z(P)=\gamma_3(P)$.
Thus $d=x^{pi}$ and $e=x^{pj}$ for some $i,j$; collection into
$x^\alpha y^\beta c^\gamma$ gives the bound $p^4$.
For $(u,v)=(1,2)$ interchange $x,y$ in this central-power argument.

In each case these full relations, together with $b^2=y$ and
$(xb)^2=1$, present $G$ by the same surjection and order bound.
Under $a\mapsto a^{-1}$, $b\mapsto b^{-1}$ we have
$x\mapsto x^{-1}$, $y\mapsto y^{-1}$.
In a class-three group their new commutator is $c$ times an element
of the central group $\gen{d,e}$, so still has $p$th power one.
The new triple commutators are $d^{-1},e^{-1}$.
Consequently the linear relation in the $(1,1)$ case is preserved,
and in the other cases the two equations with $x^p$ or $y^p$
are both inverted. Nilpotency class at most three is preserved under
any substitution, as are the orders and $(ab)^2=1$.
This proves the required involutory automorphism for every odd $p$,
including $p=3$, without appealing to a finite catalogue.
\end{proof}

